% ----------------------
\subsection{Section 23 (April 1830)}
% ----------------------
	\label{DC23}
	
	% -----
	\subsubsection{Note on the JSP versions of this section}
	% -----
		\label{DC23-JSP}
	
		Section 23 is actually a collection of five different revelations given to five different people, likely on the same occasion. The source for the JSP is RB1. That source retains the original divisions of the revelations into separate sections. Here are the recipients of the revelations, along with the verses they received:
		
			\begin{compactitem}
				\item Oliver Cowdery (\hyperref[DC23-1]{verses 1--2})
					\begin{compactitem}
						\item \hyperref[EMS-1830-JS-APR-1830-A]{RB1 account}
					\end{compactitem}
				\item Hyrum Smith (\hyperref[DC23-3]{verse 3})
					\begin{compactitem}
						\item \hyperref[EMS-1830-JS-APR-1830-B]{RB1 account}
					\end{compactitem}
				\item Samuel Smith (\hyperref[DC23-4]{verse 4})
					\begin{compactitem}
						\item \hyperref[EMS-1830-JS-APR-1830-C]{RB1 account}
					\end{compactitem}
				\item Joseph Smith Sr. (\hyperref[DC23-5]{verse 5})
					\begin{compactitem}
						\item \hyperref[EMS-1830-JS-APR-1830-D]{RB1 account}
					\end{compactitem}
				\item Joseph Knight Sr. (\hyperref[DC23-6]{verses 6--7})
					\begin{compactitem}
						\item \hyperref[EMS-1830-JS-APR-1830-E]{RB1 account}
					\end{compactitem}
			\end{compactitem}
	
	% -----
	\subsubsection{Historical Background}
	% -----
		\label{DC23-HB}
		
		% --
		\paragraph{To Oliver Cowdery, verses 1--2}
		% --
		
			%
			\subparagraph{JSP}
			%
		
				``Each of the following five revelations, which were dictated soon after the organization of the Church of Christ on 6 April 1830, addressed one of JS's family members or a close associate who desired to know the Lord's will concerning himself. All five texts include similar content and phrasing, and JS likely dictated them one after the other. John Whitmer recorded them separately in Revelation Book 1 and assigned the date `AD 1830' to each one. Though the editors of the Book of Commandments printed the revelations separately and gave each the date of 6 April 1830, that date appears to be in error and was dropped two years later, in 1835, when the Doctrine and Covenants combined the texts into a single document with a general `April, 1830' date. JS's history and other sources suggest that the revelations date between the 6 April organization and an 11 April meeting, both of which took place in Fayette Township, New York.''
		
		% --
		\paragraph{To Hyrum Smith, verse 3}
		% --
		
			%
			\subparagraph{JSP}
			%
				
				``Soon after the 6 April 1830 organizational meeting, JS dictated revelations to five individuals who were `anxious to know of the Lord what might be their respective duties' now that the church was formally established. This revelation was directed to one of the five, JS's older brother Hyrum Smith. Almost a year earlier, Hyrum visited Harmony, Pennsylvania, while JS was translating the Book of Mormon. During that visit JS announced a revelation that instructed Hyrum, `Seek not to declare my word, but first seek to obtain my word, and then shall your tongues be loosed; then, if you desire you shall have my Spirit, and my word: Yea, the power of God unto the convincing of men: but now hold your peace.' During the next year, Hyrum testified as one of the Eight Witnesses to the gold plates, safeguarded the Book of Mormon manuscript during printing, and proofread the printer's copy of the Book of Mormon.''

		% --
		\paragraph{To Samuel Smith, verse 4}
		% --
		
			%
			\subparagraph{JSP}
			%
				
				``JS dictated this revelation for his brother Samuel Harrison Smith shortly after the formal organization of the Church of Christ on 6 April 1830. It was one of five revelations that came after Samuel Smith, Hyrum Smith, Oliver Cowdery, Joseph Smith Sr., and Joseph Knight Sr. showed an eagerness to learn their respective responsibilities as members of the newly founded church. In May 1829, Samuel was the third believer to be baptized, and he later held and examined the plates and testified of their existence.''
		
		% --
		\paragraph{To Joseph Smith Sr., verse 5}
		% --
		
			%
			\subparagraph{JSP}
			%
			
				``Joseph Smith Sr. was one of JS's earliest supporters and confidants. He was the first person JS told about the visitation of Moroni and later became one of the witnesses of the gold plates. As early as February 1829, Joseph Sr. asked JS what the Lord required of him and subsequently became one of the first recipients of JS's recorded revelations. This, then, was the second revelation directed to Joseph Sr. It was one of five similarly worded revelations directed to individuals who desired to know the Lord's will.
				
				``JS's history states that Joseph Smith Sr. `came forward shortly after' the 6 April 1830 meeting and was baptized. Lucy Mack Smith recalled, `Joseph stood on the shore when his father came out of the water and as he took him by the hand he cried out Oh! my God I have lived to see my father baptized into the true church of jesus christ and he covered his face in his fathers bosom and wept aloud for joy.'''
		
		% --
		\paragraph{To Joseph Knight Sr., verses 6--7}
		% --
		
			%
			\subparagraph{JSP}
			%
			
				``Joseph Knight Sr., of Colesville, New York, probably met JS in 1826, when he hired JS as a laborer. Knight proved to be a loyal friend to JS, frequently providing money, food, supplies, and other assistance. He was particularly helpful during April and May of 1829, when JS and Oliver Cowdery were translating the gold plates in Harmony, Pennsylvania. Though an early supporter and believer, Knight was not baptized a member of the Church of Christ until June 1830. Knight was among those anxiously seeking to know their duty in light of the organization of the new church on 6 April 1830; JS subsequently dictated this revelation. Implying that Knight was hesitant to pray publicly, the revelation included the admonition to `pray vocally before the World.'''
			
	% -----
	\subsubsection{Versions}
	% -----
		\label{DC23-V}
		
		See the \href{https://drive.google.com/file/d/1goAvNz8jS4R8slYfMG_db55Ty2z_vmgK/view?usp=sharing}{sections comparison} document.
	
		\begin{compactdesc}
			\item[JSP] See \hyperref[DC23-JSP]{above} for the list.
			\item[BC] Chapters 17--21, 43--45
			\item[DC 1835] Section 45, 176--177
			\item[TS] 4:1 (15 November 1842), 12
			\item[MS] 4:8 (December 1843), 116
			\item[DC 1844] Section 45, 264--265
		\end{compactdesc}

	% -----
	\subsubsection{Section 23:1} % *DC23-1 
	% -----
		\label{DC23-1}

		BEHOLD, I speak unto you, Oliver, a few words. Behold, thou art blessed, and art under no condemnation. But beware of pride, lest thou shouldst enter into temptation.

		% \begin{framed}
		% \scriptsize
			% \begin{compactdesc}
				% \item[JSP] 
				% % \item[BC] 
				% % \item[]
			% \end{compactdesc}
		% \end{framed}

	% -----
	\subsubsection{Section 23:2} % *DC23-2 
	% -----
		\label{DC23-2}

		Make known thy calling unto the church, and also before the world, and thy heart shall be opened to preach the truth from henceforth and forever. Amen.

		% \begin{framed}
		% \scriptsize
			% \begin{compactdesc}
				% \item[JSP] 
				% % \item[BC] 
				% % \item[]
			% \end{compactdesc}
		% \end{framed}

	% -----
	\subsubsection{Section 23:3} % *DC23-3 
	% -----
		\label{DC23-3}

		Behold, I speak unto you, Hyrum, a few words; for thou also art under no condemnation, and thy heart is opened, and thy tongue loosed;%
			\footnote{%
				% \label{}%
				Compare the instruction given to Hyrum at \hyperref[DC11-21]{DC 11:21}.
				}
		and thy calling is to exhortation, and to strengthen the church continually. Wherefore thy duty is unto the church forever, and this because of thy family. Amen.

		% \begin{framed}
		% \scriptsize
			% \begin{compactdesc}
				% \item[JSP] 
				% % \item[BC] 
				% % \item[]
			% \end{compactdesc}
		% \end{framed}

	% -----
	\subsubsection{Section 23:4} % *DC23-4 
	% -----
		\label{DC23-4}

		Behold, I speak a few words unto you, Samuel; for thou also art under no condemnation, and thy calling is to exhortation, and to strengthen the church; and thou art not as yet called to preach before the world. Amen.

		% \begin{framed}
		% \scriptsize
			% \begin{compactdesc}
				% \item[JSP] 
				% % \item[BC] 
				% % \item[]
			% \end{compactdesc}
		% \end{framed}

	% -----
	\subsubsection{Section 23:5} % *DC23-5 
	% -----
		\label{DC23-5}

		Behold, I speak a few words unto you, Joseph; for thou also art under no condemnation, and thy calling also is to exhortation, and to strengthen the church; and this is thy duty from henceforth and forever. Amen.

		% \begin{framed}
		% \scriptsize
			% \begin{compactdesc}
				% \item[JSP] 
				% % \item[BC] 
				% % \item[]
			% \end{compactdesc}
		% \end{framed}

	% -----
	\subsubsection{Section 23:6} % *DC23-6 
	% -----
		\label{DC23-6}

		Behold, I manifest unto you, Joseph Knight, by these words, that you must take up your cross,%
			\footnote{%
				% \label{}%
				See \hyperref[MATTHEW-16-24]{Matthew 16:24} and notes there.
				}
		in the which you must pray vocally before the world as well as in secret, and in your family, and among your friends, and in all places.

		% \begin{framed}
		% \scriptsize
			% \begin{compactdesc}
				% \item[JSP] 
				% % \item[BC] 
				% % \item[]
			% \end{compactdesc}
		% \end{framed}

	% -----
	\subsubsection{Section 23:7} % *DC23-7 
	% -----
		\label{DC23-7}

		And, behold, it is your duty to unite with the true church, and give your language to exhortation continually, that you may receive the reward of the laborer. Amen.

		% \begin{framed}
		% \scriptsize
			% \begin{compactdesc}
				% \item[JSP] 
				% % \item[BC] 
				% % \item[]
			% \end{compactdesc}
		% \end{framed}